\exercisesection{}

¶ 1) 对 R 的每个紧区间 K，用 \(\mathcal{C}(K)\) 表示 K 上连续函数的实向量空间，赋以一致收敛拓扑；用 \(\mathcal{D}(K)\) 表示 R 上无穷可微且在 K 之外为零的实函数所成的向量空间。给 \(\mathcal{D}(K)\) 赋以使映射 \(f \mapsto D^{p}f|K\)（对每个正整数 \(p\)）都连续的最粗拓扑（D 是求导算子）。置 \(\mathcal{D} = \bigcup_{K} \mathcal{D}(K)\)，并给这个空间赋以这样的局部凸拓扑：它是子空间

\(\mathcal{D}(\mathbf{K})\) 的诸拓扑的直极限。

\begin{enumerate}
\def\labelenumi{\alph{enumi})}
\tightlist
\item
  对每个整数 \(p \geqslant 0\) 与每个函数 \(f \in \mathcal{D}(\mathbf{K})\)，置
\end{enumerate}

\[
\mathrm{Q} _ {p} (f) = \int_ {\mathrm{K}} \left(\mathrm{D} ^ {p} f (x)\right) ^ {2} d x.
\]

证明 \(\mathbf{Q}_p\) 是 \(\mathcal{D}(\mathbf{K})\) 上的一个正二次型，范数序列 \(\mathbf{Q}_p^{1/2}\) 定义了 \(\mathcal{D}(\mathbf{K})\) 的拓扑，且 \(\operatorname{Tr}(\mathbf{Q}_{p+1}/\mathbf{Q}_p)\)（对每个 \(p \geqslant 0\)）有限。（对每个实数 \(t\)，用 \(\mathrm{I}_t\) 表示区间 \([t, +\infty[\) 的特征函数；利用公式 \(\mathrm{D}^p f(t) = \int_{\mathbf{K}} \mathrm{D}^{p+1} f \cdot \mathrm{I}_t dx\) 与 Bessel 不等式，证明 \(\sum_{i=1}^{n} \mathbf{Q}_p(f_i) \leqslant l^2\) 对每个有限序列 \(f_1, \ldots, f_n\)（由属于 \(\mathcal{D}(\mathbf{K})\) 的函数组成，对 \(\mathbf{Q}_{p+1}\) 标准正交）成立；数 \(l\) 是区间 K 的长度。）由此推出 \(\mathcal{D}(\mathbf{K})\) 是一个核空间。

\begin{enumerate}
\def\labelenumi{\alph{enumi})}
\setcounter{enumi}{1}
\tightlist
\item
  证明空间 \(\mathcal{D}\) 是核的。（先确立 \(\mathcal{D}\) 中一个非零函数的存在性，并由此推出函数 \(h \geqslant 0\)（\(\mathcal{D}\) 中的）的存在性，使得 \(\sum_{n \in \mathbf{Z}} h(x - n) = 1\)（对每个 \(x \in \mathbf{R}\)）。设 \(K\) 是 \(\mathbf{R}\) 的一个紧区间，使得 \(h\) 在
\end{enumerate}

K 之外为零；对每个整数 \(n\)，置 \(h_n(x) = h(x - n)\) 与 \(\mathrm{K}_n = \mathrm{K} + n\)。设 V 是 \(\mathcal{D}\) 中 0 的一个凸邻域；存在正整数 \(p_n\)，使得每个函数 \(f \in \mathcal{D}(\mathrm{K}_n)\)（满足 \(\mathrm{Q}_{p_n}(f) \leqslant 1\)）都属于 V。用下列公式定义 \(\mathcal{D}\) 上的连续二次型 Q 与 R：\(\mathrm{Q}(f) = \sum_{n \in \mathbf{Z}} 2^{2n} \mathrm{Q}_{p_n}(f \cdot h_n)\) 与 \(\mathrm{R}(f) = \sum_{n \in \mathbf{Z}} 2^{3n} \mathrm{Q}_{p_n + 1}(f \cdot h_n)\)；证明 V 包含由 \(f \in \mathcal{D}\)（满足 \(\mathrm{Q}(f) \leqslant 1\)）组成的集合，且 \(\operatorname{Tr}(\mathrm{R}/\mathrm{Q})\) 有限。）

\(^{*}\) c) 把上述结果推广到 \(R^{n}\) 上具有紧支撑的无穷可微函数。\(^{*}\)

